% Partie : théorie de la démonstration
% Chapitre : normalisation
% Section : théorème de normalisation

\documentclass[../../../include/open-logic-section]{subfiles}

\begin{document}

\olfileid[fr]{pt}{nor}{thm}

\olsection{Le théorème de normalisation}

Une !!{proof} est en \emph{forme normale} si elle
ne comporte aucun segment de coupure. Avec les
conversions définies précédemment, cela signifie
qu'aucune conversion par permutation ou par
réduction ne lui est applicable.
\emph{Normaliser} une preuve consiste à la
transformer en une preuve normale par ces
conversions, jusqu'à disparition de toutes les
coupures. Le \emph{théorème de normalisation}
affirme qu'une telle transformation est
toujours possible.

\begin{thm}\ollabel{thm:normalization}
Toute !!{proof}~$\delta$ de \Log{N1i} peut
être normalisée, c'est-à-dire transformée
en une !!{proof}~$\delta'$ sans coupure,
de même conclusion, dont les hypothèses
ouvertes sont parmi celles de la preuve
initiale.
\end{thm}

\begin{proof}
Nous raisonnons par récurrence lexicographique
sur le rang de coupure et la longueur de
coupure de~$\delta$. L'hypothèse de récurrence
porte sur toutes les preuves de rang plus
petit, ou de même rang et de longueur de
coupure plus petite.

Si la longueur de coupure est~$0$, la preuve
n'a aucune coupure : elle est déjà normale.
Supposons donc sa longueur de coupure~$>0$.

Choisissons une coupure maximale telle
qu'aucune autre coupure maximale ne se
termine au-dessus de sa première occurrence,
ni au-dessus d'une autre prémisse de
l'élimination dont sa dernière occurrence
est la prémisse majeure, et qu'aucune autre
coupure maximale ne contienne une telle
prémisse voisine.
On obtient une telle coupure en choisissant, parmi les coupures
maximales qui n'en ont aucune au-dessus, celle dont la première
occurrence est le plus à droite, les prémisses majeures étant
placées à gauche. Toute coupure maximale au-dessus d'une prémisse
voisine, ou qui la contiendrait, commencerait dans une branche à
droite. En remontant au besoin à une coupure maximale supérieure
dans cette branche, on contredirait le choix le plus à droite.


Si le segment choisi a une longueur~$>1$,
on applique la conversion par permutation
correspondante. S'il a une longueur~$1$,
on applique la conversion par réduction
correspondante. Les conditions du choix
empêchent de dupliquer une coupure
maximale ou d'allonger un segment maximal
passant par une prémisse voisine.
Les segments visés sont raccourcis ou
supprimés ; les nouvelles coupures
créées par une réduction ont un rang
inférieur. On obtient donc une preuve
de rang plus petit, ou de même rang et
de longueur de coupure plus petite.
Après les renommages nécessaires pour
éviter la capture et rétablir la régularité,
l'hypothèse de récurrence s'applique.
Les conversions préservent la conclusion
et n'ajoutent aucune hypothèse ouverte.
\end{proof}

La démonstration donne une stratégie
particulière : elle traite toujours une
coupure maximale choisie de cette façon.
Elle établit l'existence d'une suite
terminante de conversions, appelée
\emph{normalisation faible}.
La \emph{normalisation forte} affirme que toute suite maximale
de conversions est finie et s'achève en forme normale.
Sa démonstration exige un argument de terminaison indépendant
du choix de la coupure. La mesure ci-dessus décroît sous
les conditions de sélection de cette preuve.






Les définitions de segment et de coupure
se transposent à \Log{N2i}, en considérant
les occurrences de la formule à droite
de chaque séquent. Les conversions de
\Log{N1i} se transposent elles aussi :
on recalcule les contextes d'hypothèses
ouvertes après greffage ou suppression
d'une branche. Les décharges et les
conditions sur les constantes propres
sont les mêmes. Le théorème vaut donc
également pour~\Log{N2i}.

\begin{thm}\ollabel{thm:normalization-N2i}
Toute !!{proof}~$\delta$ de \Log{N2i} peut
être normalisée, c'est-à-dire transformée
en une !!{proof}~$\delta'$ sans coupure,
de même formule finale, dont le contexte
d'hypothèses ouvertes est inclus dans
celui de la preuve initiale.
\end{thm}

Pour le choix du segment maximal dans la preuve de Prawitz, voir
la section~5 de l'article d'Ernst Zimmermann,
\href{https://doi.org/10.1007/s10849-021-09329-8}{\emph{Natural Deduction Bottom Up}} (2021).

\end{document}















